\documentclass[12pt]{article}
\usepackage[utf8]{inputenc}
\usepackage[T1]{fontenc}
\usepackage{hyperref}
\title{Jaitmatang Language Resources for K-12 Curriculum Development}
\author{Compiled by Claude Sonnet 5.5 (claude-sonnet-5-5) for Andr\'as Kornai}
\date{}
\begin{document}
\maketitle

\section{Introduction}
Jaitmatang (ISO 639-3 xjt; also spelled Yaithmathang or Jaithmathang) was spoken in north-east Victoria, Australia, around the headwaters of the Mitta Mitta and Tambo rivers. It is described as extinct. Sources say its relationship to neighbouring Ngarigo and Dhudhuroa is debated, and one view treats it as a dialect of Ngarigo. The search results show academic reconstruction and community cultural-reconnection work. They show no Jaitmatang-specific language programme and no teaching materials.

\section{Primary Educational Resources}
No primers, textbooks or school programmes for Jaitmatang were found. The only education-related item is the Victorian VCE study design for Aboriginal languages of Victoria, which covers revival and reclamation in general terms.

\section{Linguistic and Reference Materials}
\begin{itemize}
\item Wikipedia overview of Yaithmathang: \url{https://en.wikipedia.org/wiki/Yaithmathang}
\item ``Dhudhuroa and Yaithmathang languages and social groups in north-east Victoria -- a reconstruction'' (ANU Press chapter): \url{https://press-files.anu.edu.au/downloads/press/p74631/html/ch09.html?page=12}
\item High Country Cosmos, A--Z culture entry on Jaithmathang: \url{https://www.hgcosmos.org/assets/culture/AZ/J/Jaithmathang.pdf}
\item Sydney University Press, ``Increasing the accessibility of information on the indigenous languages of Victoria'': \url{https://open.sydneyuniversitypress.com.au/9781920899554/rl-increasing-the-accessibility-of-information-on-the-indigenous-languages-of-victoria.html}
\end{itemize}

\section{Religious and Community Resources}
\begin{itemize}
\item The Mountain Journal, ``Jaithmathang senior elders reconnect with their original country'' (2021): \url{https://themountainjournal.com/2021/06/29/jaithmathang-senior-elders-reconnect-with-their-original-country/}
\item North East Catchment Management Authority news item on the same project: \url{https://www.necma.vic.gov.au/News-Events/News/ArtMID/431/ArticleID/1629/Jaithmathang-senior-elders-reconnect-with-their-ancestral-MPP-and-original-country}
\end{itemize}
No scripture translations were found.

\section{Digital and Audio Resources}
None were found for Jaitmatang.

\section{Curriculum Development Context}
The language is extinct, so any curriculum would be reclamation work. It would rest on the historical word lists discussed in the ANU Press reconstruction chapter and be led by Jaithmathang community members. Victorian structures for language revival exist but were not shown to cover Jaitmatang. They are the Victorian Aboriginal Corporation for Languages and the VCAA study design for Aboriginal languages of Victoria: \url{https://vcaa.vic.edu.au/sites/default/files/2026-06/2023AboriginalLanguagesVictoriaSD.docx}. Very little material exists, and the dossier does not claim otherwise.

\end{document}